% TOP_PROBLEM: {"id": 2935, "title": "Kirby Problem 4.59", "category_id": 7, "category": {"id": 7, "name": "topology", "display_name": "Topology", "description": "Properties preserved under continuous deformations.", "slug": "topology", "order_index": 7, "created_at": "2026-07-31T15:26:25.671Z"}, "status": "open", "problem_number": "KP-4.59", "set_id": 11, "statusQualification": "Integral coefficients and consistent oriented cobordism/homeomorphism conventions are explicit."}
% FIRST_POSTED: 2026-10-01
% ENTRY_KIND: statement_only
% UnsolvedMath ID 2935; source catalogue code KP-4.59
% !TeX program = lualatex
% Definition and sources only; this file is not an LLM solution attempt.
\documentclass[11pt,a4paper]{article}
\usepackage[margin=25mm]{geometry}
\usepackage{fontspec}
\setmainfont{lmroman10-regular.otf}[BoldFont=lmroman10-bold.otf,
 ItalicFont=lmroman10-italic.otf,BoldItalicFont=lmroman10-bolditalic.otf]
\usepackage{amsmath,amssymb,mathtools}
\usepackage{unicode-math}
\setmathfont{latinmodern-math.otf}
\AtBeginDocument{\let\setminus\smallsetminus}
\usepackage{xurl,hyperref}
\hypersetup{colorlinks=true,urlcolor=blue}
\setcounter{secnumdepth}{0}
\setlength{\parindent}{0pt}
\setlength{\parskip}{5pt}
\setlength{\emergencystretch}{3em}
\begin{document}
\section*{Kirby Problem 4.59}\label{2935}
{\small UnsolvedMath ID 2935 · KP-4.59 · Topology · Kirby's Problems in Low-Dimensional Topology}
\subsection{Definitions and mathematical statement}
For integers $p\ge1$ and $q$ coprime to $p$, let $L(p,q)$ be the
lens space obtained from $S^3\subset\mathbb C^2$ by the free cyclic
action generated by
$(z_1,z_2)\mapsto(e^{2\pi i/p}z_1,e^{2\pi iq/p}z_2)$.
Use the quotient orientation. It has first integral homology
$\mathbb Z/p\mathbb Z$.

Two oriented closed $3$-manifolds $Y_0,Y_1$ are topologically
integral-homology cobordant if there is a compact oriented topological
$4$-manifold $W$ with $\partial W=(-Y_0)\sqcup Y_1$ such that
each boundary inclusion induces an isomorphism
\[
 H_i(Y_j;\mathbb Z)\longrightarrow H_i(W;\mathbb Z)
          \quad\text{for every }i,\ j=0,1.
\]
No smooth structure on $W$ is required.

Are two lens spaces homology cobordant in this sense if and only
if they are orientation-preservingly homeomorphic? The homeomorphism
implies the cobordism direction by taking a product. The difficult
direction asks whether an arbitrary topological homology cobordism
forces that homeomorphism.

The coefficients are integers, not rationals: rational homology
cobordism is a weaker relation. The hypotheses immediately force the
two lens spaces to have the same $p$, since their first homology
groups are isomorphic, but the gluing parameter $q$ still matters.
Orientations are compared consistently; reversing the chosen orientation
of one lens space changes the oriented instance of the question.
\subsection{Short English statement}
Are lens spaces topologically integral-homology cobordant exactly when they are homeomorphic, with the corresponding orientation convention?
\subsection{Sources}
\begin{itemize}
\item[S1] ULAM.AI and the UnsolvedMath Contributors, supplied problems.json, record 2935 (KP-4.59), Kirby's Problems in Low-Dimensional Topology. Catalogue identity and source transcription; CC BY 4.0.
\url{https://huggingface.co/datasets/ulamai/UnsolvedMath}.
\item[S2] K3: A New Problem List in Low-Dimensional Topology, Problem 4.59. Expanded formulation of the supplied catalogue statement.
\url{https://math.berkeley.edu/sites/default/files/surv-295-ruberman-watermarked-author-pdf.pdf}.
\end{itemize}
\end{document}
